\subsection{Unconditional resource bounds}\label{sec:resources}

A surviving representation must reach a completed check before the
operation cap. Recall the three ambient counts from the prefix
analysis: \(G\) counts low-filtered raw child pairs, \(J_j\) counts
low-filtered raw four-leaf tuples on side \(j\), and \(U_j\) counts
distinct valid weight--top-mask entries passing only the parent low
filter. None uses \(q\), a prefix restriction, or a separating family.
The first two retain all occurrence multiplicities, including tuples
with a local overlap. The last counts only valid entries and ignores
the child filters. We shall bound these counts without conditioning on
the rare survival event.

For their first moments, condition on a passing mixer setup, signs,
guesses, layout and allowed moduli \(p_i\), leaving the low residues
\(\rho_0,\rho_L,\rho_R\) independent and uniform. Each child's required
residue is marginally uniform modulo \(p_j\), including child~2 because
\(\rho_j\) is independent of \(\rho_0\). Each fixed four-leaf tuple
passes its parent filter and first-child filter with probability
\(1/(p_0p_j)\); divisibility \(p_j\mid p_0\) then gives the second-child
filter. Before its parent filter, the number of distinct valid entries
on side \(j\) is at most
\[
 2^{r_j}\,2^{H_{0j}m}\,2^{m-d_j}.
\]
Here the three factors count background choices, possible top masks,
and distinct actual weights from disjoint use of \(M_j\). The last
factor is \(N_j\), also for the signed weights. Each entry passes the
parent filter with probability \(1/p_0\).

Define the analysis scales
\begin{equation}\label{res:scales}
 \begin{aligned}
 \Gamma&=\frac n4+\left(-\frac54+\frac32h_*\right)m+f_{\max},\\
 d_*&=\min(d_L,d_R),\\
 \Lambda&=\frac n2+(\overline H_0-1)m-d_*+f_0.
 \end{aligned}
\end{equation}
By \eqref{eq:allocation}, \eqref{eq:beta}, and
\eqref{eq:filter-bounds}, the conditional moments satisfy
\begin{equation}\label{res:moments}
 \mathbb EG\le u^{26}2^\Gamma,\qquad
 \mathbb E(J_L+J_R)\le u^{51}2^{n/2},\qquad
 \mathbb EU_j\le u^{10}2^\Lambda.
\end{equation}
For the first bound, use \(\beta_j-L_j\le\Gamma\), the
\(u^{20}\) factor for raw pairs, and the \(u^5\) loss in the
modulus lower bound. For the second, raw four-leaf tuple counts are at
most \(u^{40}2^{2\beta_j}\), and their exponent after filtering obeys
\begin{equation}\label{res:tuple-exponent}
 \begin{aligned}
 2\beta_j-L_0-L_j
 &\le \frac n2+\left(-\frac52+3h_*\right)m+2\tau\\
 &\le \frac n2.
 \end{aligned}
\end{equation}
The final inequality holds for large \(n\), since the excess is at
most \(-.0658m+.0156n\), and \(m=.24n+O(1)\). For the third
bound, the exponential part of the preceding universe count divided by
the parent modulus is
\[
 \frac r2+\overline H_0m+m-d_j-L_0
 =\frac n2+(\overline H_0-1)m-d_j+f_0
 \le\Lambda;
\]
the allocation's logarithmic error and the modulus loss fit in the
factor \(u^{10}\).

Let \(\mathcal E\) be the resource event
\begin{equation}\label{eq:resource-event}
 G\le u^{200}2^\Gamma,\qquad
 J_L+J_R\le u^{200}2^{n/2},\qquad
 U_L,U_R\le u^{200}2^\Lambda.
\end{equation}
The scales in \eqref{res:scales} are fixed under the conditioning used
in \eqref{res:moments}. Markov's inequality and a union bound, followed
by averaging over that conditioning, give
\begin{equation}\label{res:joint-failure}
 \Pr[\text{spread test passes and }\mathcal E^c]
 \le u^{-174}+u^{-149}+2u^{-190}
 <\frac1{10}u^{-100}.
\end{equation}
This is a joint failure estimate, not an estimate conditioned on the
rare event in Lemma~\ref{lem:representation}. Consequently that lemma's
surviving-representation event and \(\mathcal E\) hold together with
probability at least \(.9u^{-100}\) in the low-mean YES case.

Three exponent relations control the subsequent costs:
\begin{equation}\label{res:exponent-relations}
 \xi+\Gamma\le\frac n2,\qquad
 \Lambda+(1-h_*)m\le\frac n2,\qquad
 \xi-(\Lambda-.199n)\ge .0015n.
\end{equation}
The first relation pays for scanning the low-filtered children over
roughly \(2^\xi\) tasks. The second pays for the repeated separator
incidences of distinct records. The third makes the expected number
of different-weight records in a solution bucket smaller than its
capacity \(C\). These relations hold for large \(n\) on every passing
setup. The first follows
from \(.032-1.25+1.5h_*<0\), with cancellation of \(f_{\max}\).
For the second, use \(\overline H_0\le h_*\) and
\(f_0\le d_*\), since \(M_0\) has least deficiency. For the
third, direct subtraction gives the lower bound
\[
 -.051n+(1.032-h_*)m+d_*-f_0-f_{\max}.
\]
Here \(f_0\le d_*/2\), while the spread test implies
\(f_{\max}\le d_*/2+.00025n\). Since
\(1.032-h_*>.2206\) and \(m=.24n+O(1)\), this lower bound is
at least \(.0015n\). The first two relations preserve the exact
time exponent \(1/2\); no positive exponential slack is added there.

\subsection{The extra prime and the solution bucket}

Condition on a surviving representation and \(\mathcal E\), exposing
all choices so far but still not \(q\). In its loop iteration, the
expected number of entries in side \(j\)'s universe with weight
different from \(x_j\) but congruent to it modulo \(P=p_0q\) is
at most
\begin{equation}\label{res:bucket-expectation}
 U_j u^3 2^{-\xi}
 \le u^{203}2^{.199n-.0015n}.
\end{equation}
The nonzero difference must be divisible by the independent prime
\(q\), so the prime collision bound applies. No coprimality of
\(q\) and \(p_0\) is required. Since
\(C\ge 2^{.199n-1}\), the right side of
\eqref{res:bucket-expectation} divided by \(C/2\) is at most
\[
 2^{(203/10^9-.0015)n+2}
\]
on the main branch. Markov's inequality shows, with conditional
probability at least \(.9\) for large \(n\), that both sides have
at most \(C/2\) such different-weight entries.

Equal-weight entries need a separate argument. On L, at most
\(\binom m{k_0}\le 2^{h_*m}\le 2^{.194736n}<C/2\) distinct
entries can have weight \(x_L\), since the smaller top mask has
size at most \(m/4\). Thus the entire solution L bucket fits. On R,
follow the prefix of the solution top mask. At a full prefix there is
at most one entry of weight \(x_R\), by deduplication. Together
with the different-weight bound, this full-prefix bucket fits, so the
solution's prefix path terminates in a completed nonoverflow task.
In every task compatible with the solution masks, its surviving leaf
choices generate the corresponding parent entries. Hence the solution
pair reaches a weighted disjointness check, unless an earlier YES or
an operation or sampling abort has already occurred.

\subsection{Applying the three prefix bounds}

Consider the complete hypothetical parent schedule obtained by ignoring
YES answers and the operation cap, while retaining the actual overflow
rules. Its generation uses only the parent choices and \(q\). Thus it
can be fixed without examining any separating family. The three
bounds in Lemma~\ref{lem:prefix} now have their required numerical
inputs; we apply each to the quantity it controls.

First, the number of parent-generation tasks is at most
\begin{equation}\label{res:task-count}
 T\le 2q+2(m+1)U_R/C\le u^{203}2^\xi.
\end{equation}
Indeed, \(q\le2^\xi\), \(C\ge2^{.199n-1}\), and the last inequality in
\eqref{res:exponent-relations} gives
\(U_R/C\le2u^{200}2^{\xi-.0015n}\) on \(\mathcal E\).
The main-branch guard absorbs the remaining polynomial factor.

Second, the concrete stopped-stream implementation supplies
\begin{equation}\label{res:nonmatch-cost}
 H=u^8(2^{.2n}+G),\qquad K=u^8
\end{equation}
in the task-cost formula of Lemma~\ref{lem:prefix}. The term \(H\)
pays for regenerating the leaves, prefix tests before streaming,
sorting, stream initialization and traversal, the bounded tie buffer,
and dictionary initialization. The term \(K\) pays for each potential
matching tuple, including replay, invalid local overlaps and repeated
entries. Consequently the full parent-generation work is at most
\[
 \begin{aligned}
 H T+K\bigl(J_L+(m+1)J_R\bigr)
 &\le u^{211}2^\xi(2^{.2n}+u^{200}2^\Gamma)
       +u^8\bigl(J_L+(m+1)J_R\bigr)\\
 &\le u^{420}2^{n/2}
 \end{aligned}
\]
for sufficiently large main-branch inputs. Here \(\xi+.2n<n/2\) follows
from \(\xi\le.25768n\), whereas \(\xi+\Gamma\le n/2\) is the first
exact relation in \eqref{res:exponent-relations}. The raw-match term
is charged by prefix depth; it is never multiplied by \(T\).

Third, write \(N_{\mathrm{check}}\) for the total array length supplied
to all final checks, counting every reuse of a left array. The final
bound of Lemma~\ref{lem:prefix} gives
\begin{equation}\label{res:check-inputs}
 N_{\mathrm{check}}\le U_L+(2m+3)U_R.
\end{equation}
Completed right chunks pay for their own entries, and overflow entries
pay for repeated left arrays. This is why the final checks use the
distinct-entry counts \(U_j\), even though generation required the raw
counts \(J_j\).

\subsection{Independent families and bounded-trial success}

Fix all choices through \(q\), on the joint survival, resource and
bucket events. Choose deterministically the check on the completed
solution-prefix path and the solution pair in it. The complete
hypothetical check list, its layout parameters, this designated check
and its pair are all functions of those nonfamily choices. Conditional
on the layout, the program sampled the four indexed families from the
product law specified in Section~\ref{sec:separator}, independently
of every parent and filter choice. Therefore that same product law
holds after conditioning jointly on this entire list, its parameters,
and the designated pair. This fact, rather than delayed exposure by
itself, verifies the hypothesis of Corollary~\ref{cor:separator-reuse}.

On \(\mathcal E\), \eqref{res:check-inputs} and that corollary give
conditional expected total check work at most
\[
 N_{\mathrm{check}}u^{22}2^{(1-h_*)m}
 \le u^{230}2^{\Lambda+(1-h_*)m}
 \le u^{230}2^{n/2}.
\]
The second exact relation in \eqref{res:exponent-relations} is used
here. Markov's inequality bounds the chance of exceeding
\(u^{240}2^{n/2}\) operations by \(u^{-10}\). The designated pair has
coverage probability at least \(1-4e^{-10}\). A union bound gives
both coverage and this work bound with conditional probability at
least \(1-4e^{-10}-u^{-10}\); these events need not be independent.
The hypothetical list is an analysis device and is never stored.

\begin{lemma}\label{lem:resources}
On a YES input with \(\mu<.00765n\), a representation trial with
the bounded samplers and operation cap succeeds with probability at
least \(u^{-110}\), for sufficiently large main-branch inputs.
\end{lemma}

\begin{proof}
Lemma~\ref{lem:representation} and \eqref{res:joint-failure} give
probability at least \(.9u^{-100}\) of simultaneous survival and
resources. The independent extra prime gives the bucket event with
conditional probability at least \(.9\). Given both stages, the
family argument gives coverage and check work with conditional
probability at least \(1-4e^{-10}-u^{-10}\). These multiply as
conditional lower bounds, not as an assertion that the events are
independent. On their intersection, parent generation costs at most
\(u^{420}2^{n/2}\) and all checks cost at most \(u^{240}2^{n/2}\).

The remaining setup work is smaller. Counting the three mixer
sumsets by half-list streams takes \(\poly(n)2^m\) operations.
There are constantly many prime calls, each using at most \(n^3\)
attempts, with \(\poly(n)2^{.129n}\) work per attempt. Family
sampling and incidence construction take \(\poly(n)2^{2g}\)
operations, including their bounded range-draw costs. Direct
combination generation and weight/mask formation cost at most
\(O(n^2)\) per item; comparison sorting, heap operations and interval
searches cost at most \(O(n)\) per item or advance; trie operations
cost \(O(u)\) per key. These explicit operations fit the displayed
loose polynomial powers. Including the worst permitted sampler costs,
the joint favorable execution uses fewer than
\(u^{500}2^{\lceil n/2\rceil}\) underlying word operations for large
\(n\).

Lemma~\ref{lem:samplers} couples the bounded draws to the ideal
experiment with loss \(\exp(-\Omega(n^2))\) per trial. This is
smaller than \(u^{-110}\) by any fixed factor on the main branch.
Subtracting that loss from
\[
 .81u^{-100}(1-4e^{-10}-u^{-10})
\]
leaves at least \(u^{-110}\). On this intersection the trial completes
before its cap and finds a solution. An earlier YES is already a
success.
\end{proof}
